% Heart Sutra Synoptic Critical Edition - Prototype v2
% With inline superscripts and critical apparatus
% Landscape A4

\documentclass[10pt,landscape,a4paper]{article}

% ============================================================================
% PACKAGES
% ============================================================================

\usepackage[landscape,margin=1.2cm,bottom=1.5cm]{geometry}

% Page numbers in bottom right corner
\usepackage{fancyhdr}
\pagestyle{fancy}
\fancyhf{}
\renewcommand{\headrulewidth}{0pt}
\fancyfoot[R]{\thepage}
\usepackage{fontspec}
\usepackage{xeCJK}
\usepackage{tabularx}
\usepackage{array}
\usepackage{booktabs}
\usepackage{xcolor}

% Footnotes - allow multiple references, better formatting
\usepackage[hang,flushmargin]{footmisc}
\usepackage{footnote}
\makesavenoteenv{tabularx}

% For custom footnote markers
\usepackage{hyperref}
\hypersetup{colorlinks=true,linkcolor=blue}

% ============================================================================
% FONTS
% ============================================================================

\setmainfont{Noto Serif}
\setsansfont{Noto Sans}
\setCJKmainfont{Noto Serif CJK SC}
\setCJKsansfont{Noto Sans CJK SC}

\newfontfamily\devanagarifont{Noto Serif Devanagari}
\newcommand{\deva}[1]{{\devanagarifont #1}}

\newfontfamily\tibetanfont{Noto Serif Tibetan}
\newcommand{\tib}[1]{{\tibetanfont #1}}

% ============================================================================
% CUSTOM COMMANDS
% ============================================================================

% Segment number
\newcommand{\segnum}[1]{\textbf{[#1]}}

% Romanization (smaller, different color)
\newcommand{\rom}[1]{{\small\color{gray} #1}}

% English translation
\newcommand{\engltrans}[1]{\multicolumn{3}{p{0.95\textwidth}}{\small\itshape ``#1''}}

% Apparatus note types
\newcommand{\var}[1]{{\color{red}\textsuperscript{\textbf{#1}}}}  % variant marker - red and bold
\newcommand{\lemma}[1]{\textbf{#1}}  % lemma in apparatus
\newcommand{\reading}[1]{#1}  % variant reading
\newcommand{\wit}[1]{\textsc{#1}}  % witness siglum
\newcommand{\note}[1]{\textit{#1}}  % editorial note

% Direction of dependence markers
\newcommand{\depmark}[2]{{\small\color{blue}[#1\,$\rightarrow$\,#2]}}

% ============================================================================
% DOCUMENT
% ============================================================================

\begin{document}

% Title
\begin{center}
{\LARGE\bfseries Prajñāpāramitāhṛdaya}\\[0.2em]
{\LARGE 般若波羅蜜多心經}\\[0.2em]
{\large Heart Sūtra -- Synoptic Critical Edition}\\[0.5em]
{\small Base texts: Chinese T251 (649 CE) · Sanskrit GRETIL · Tibetan Toh 21}\\[0.3em]
{\footnotesize\itshape Following Chinese compositional priority (Nattier 1992)}
\end{center}

\vspace{0.5em}

% ============================================================================
% WITNESS SIGLA KEY
% ============================================================================

\noindent\fbox{\parbox{0.98\textwidth}{
\footnotesize
\textbf{Sigla:}
\textit{Chinese:} \wit{T250}--\wit{T257} = Taishō editions; \wit{S2464} = Dunhuang.
\textit{Sanskrit:} \wit{Ja} = Hōryū-ji; \wit{Nb} = Cambridge 1485; \wit{Nk} = Cambridge 1680; \wit{Gilgit} = Gilgit fragments.
\textit{Tibetan:} \wit{Toh21} = Degé; \wit{Stok} = Stok Palace; \wit{IOL} = Dunhuang IOL J 751.
\quad\depmark{zh}{sa} = Chinese to Sanskrit dependence.
}}

\vspace{1em}

% ============================================================================
% SYNOPTIC TABLE
% ============================================================================

\noindent
\begin{tabularx}{\textwidth}{|>{\raggedright\arraybackslash}X|>{\raggedright\arraybackslash}X|>{\raggedright\arraybackslash}X|}
\hline
\textbf{Chinese (T251)} & \textbf{Sanskrit} & \textbf{Tibetan (Toh 21)} \\
\hline
\hline

% --- Segment 1: Opening ---
\segnum{1} 觀自在\var{1}菩薩，行深般若波羅蜜多時，照見五蘊皆空，度一切苦厄。
\newline\rom{Guānzìzài púsà, xíng shēn bōrě bōluómìduō shí, zhàojiàn wǔyùn jiē kōng, dù yīqiè kǔ'è.}
&
\deva{आर्यावलोकितेश्वर}\var{2}\deva{बोधिसत्त्वो गम्भीरां प्रज्ञापारमिताचर्यां चरमाणो व्यवलोकयति स्म।}\var{3}
\newline\rom{āryāvalokiteśvara-bodhisattvo gambhīrāṃ prajñāpāramitācaryāṃ caramāṇo vyavalokayati sma.}
&
\tib{འཕགས་པ་སྤྱན་རས་གཟིགས་དབང་ཕྱུག་བྱང་ཆུབ་སེམས་དཔའ་}
\newline\rom{'phags pa spyan ras gzigs dbang phyug byang chub sems dpa'...}
\\
\hline
\engltrans{The noble Avalokiteśvara Bodhisattva, while practicing the profound prajñāpāramitā, looked down and saw that the five aggregates are all empty.} \\
\hline
\hline

% --- Segment 2: Form is emptiness ---
\segnum{2} 舍利子，色不異空，空不異色；色即是空，空即是色。受、想、行、識，亦復如是。
\newline\rom{Shèlìzǐ, sè bù yì kōng, kōng bù yì sè; sè jí shì kōng, kōng jí shì sè.}
&
\deva{इह शारिपुत्र रूपं शून्यता शून्यतैव रूपं।}\var{4} \deva{रूपान्न पृथक् शून्यता शून्यताया न पृथग् रूपम्।}
\newline\rom{iha śāriputra rūpaṃ śūnyatā śūnyataiva rūpaṃ. rūpānna pṛthak śūnyatā śūnyatāyā na pṛthag rūpam.}
&
\tib{ཤཱ་རིའི་བུ་གཟུགས་སྟོང་པའོ། སྟོང་པ་ཉིད་གཟུགས་སོ།}
\newline\rom{shā ri'i bu gzugs stong pa'o. stong pa nyid gzugs so.}
\\
\hline
\engltrans{Śāriputra, form is not different from emptiness, emptiness is not different from form; form is exactly emptiness, emptiness is exactly form.} \\
\hline
\hline

% --- Segment 3: No ignorance... ---
\segnum{3} 無無明，亦無無明盡\var{5}；乃至無老死，亦無老死盡。
\newline\rom{Wú wúmíng, yì wú wúmíng jìn; nǎizhì wú lǎosǐ, yì wú lǎosǐ jìn.}
&
\deva{न विद्या नाविद्या न विद्याक्षयो}\var{6} \deva{नाविद्याक्षयो यावन्न जरामरणं न जरामरणक्षयः।}
\newline\rom{na vidyā nāvidyā na vidyākṣayo nāvidyākṣayo yāvanna jarāmaraṇaṃ na jarāmaraṇakṣayaḥ.}
&
\tib{མ་རིག་པ་མེད། མ་རིག་པ་ཟད་པ་མེད་པ་ནས་རྒ་ཤི་མེད། རྒ་ཤི་ཟད་པའི་བར་དུ་ཡང་མེད་དོ།}
\newline\rom{ma rig pa med. ma rig pa zad pa med pa nas rga shi med. rga shi zad pa'i bar du yang med do.}
\\
\hline
\engltrans{No ignorance, and no ending of ignorance; up to no old age and death, and no ending of old age and death.} \\
\hline
\hline

% --- Segment 4: Mantra ---
\segnum{4} 即說咒曰：揭諦揭諦，波羅揭諦，波羅僧揭諦，菩提薩婆訶。
\newline\rom{Jí shuō zhòu yuē: Jiēdì jiēdì, bōluó jiēdì, bōluósēng jiēdì, pútí sàpóhē.}
&
\deva{तद्यथा गते गते पारगते पारसंगते बोधि स्वाहा।}\var{7}
\newline\rom{tadyathā gate gate pāragate pārasaṃgate bodhi svāhā.}
&
\tib{ཏདྱ་ཐཱ། ག་ཏེ་ག་ཏེ། པཱ་ར་ག་ཏེ། པཱ་ར་སཾ་ག་ཏེ། བོ་དྷི་སྭཱ་ཧཱ།}
\newline\rom{tadyathā / gate gate / pāragate / pārasaṃgate / bodhi svāhā}
\\
\hline
\engltrans{The mantra is spoken thus: Gone, gone, gone beyond, gone completely beyond, awakening, svāhā!} \\
\hline

\end{tabularx}

\vspace{1.5em}

% ============================================================================
% CRITICAL APPARATUS
% ============================================================================

\noindent\textbf{Apparatus}
\footnotesize

\noindent
\textbf{1}~\lemma{觀自在}] \reading{觀世音} \wit{T250 T252}; supports Skt. \textit{-īśvara} vs. \textit{-svara}
\quad\textbf{2}~\lemma{avalokiteśvara-}] \reading{avalokita-} \wit{Ja}; \depmark{zh}{sa} Chinese supports \textit{-īśvara}
\quad\textbf{3}~\lemma{vyavalokayati sma}] unusual impf.; cf. Chinese 時

\noindent
\textbf{4}~\lemma{rūpaṃ śūnyatā}] \reading{rūpaṃ śūnyatāṃ} \wit{Ja} (scribal error)
\quad\textbf{5}~\lemma{盡}] = Skt. \textit{kṣaya} not \textit{nirodha}; key back-translation evidence
\quad\textbf{6}~\lemma{vidyākṣayo}] \depmark{zh}{sa} \textit{kṣaya} for expected \textit{nirodha} (Nattier 1992:175--8)

\noindent
\textbf{7}~\lemma{gate gate}] mantra transliterated in all traditions; 揭諦 = phonetic \textit{gate}

\vspace{0.8em}
\noindent\textbf{Notes:}
\depmark{zh}{sa} = evidence of Chinese → Sanskrit dependence.
Key variants: \textbf{2} and \textbf{6} are primary evidence for back-translation thesis (Nattier 1992).

\vspace{1em}
\begin{center}
{\small\itshape [Prototype v2 -- 4 segments with apparatus]}
\end{center}

\end{document}
